\FigureLayoutDeclare{img/2007/zhejiang_science/q4_fig1.png}{6.0cm}{54.24bp 2.64bp 60.96bp 1.92bp}

\chapter{2007年浙江卷（理）}

\section{单选题}

\begin{problem}
“\(x > 1\)” 是 “\(x^{2} > x\)” 的（\quad）

\choices
    {充分而不必要条件}
    {必要而不充分条件}
    {充分必要条件}
    {既不充分也不必要条件}
\begin{answer}
A
\end{answer}

\begin{solution}
由 \(x>1\)，得 \(x(x-1)>0\)，所以 \(x^2>x\). 但由 \(x^2>x\) 只能得到 \(x<0\) 或 \(x>1\)，例如 \(x=-1\) 时 \(x^2>x\) 成立而 \(x>1\) 不成立. 因此“\(x>1\)”是“\(x^2>x\)”的充分而不必要条件，选 A.
\end{solution}
\end{problem}

\begin{problem}
若函数 \(f(x) = 2\sin(\omega x + \varphi)\)，\(x \in \R\) （其中 \(\omega > 0\)，\(|\varphi| < \frac{\pi}{2}\)）的最小正周期是 \( \pi \)，且 \( f(0) = \sqrt{3} \)，则（\quad）

\choices
    {\(\omega = \frac{1}{2}\)，\(\varphi = \frac{\pi}{6}\)}
    {\(\omega = \frac{1}{2}\)，\(\varphi = \frac{\pi}{3}\)}
    {\(\omega = 2\)，\(\varphi = \frac{\pi}{6}\)}
    {\(\omega = 2\)，\(\varphi = \frac{\pi}{3}\)}
\begin{answer}
D
\end{answer}

\begin{solution}
函数的最小正周期为 \(\dfrac{2\pi}{\omega}\)，故 \(\dfrac{2\pi}{\omega}=\pi\)，从而 \(\omega=2\). 又 \(f(0)=2\sin\varphi=\sqrt{3}\)，且 \(\lvert\varphi\rvert<\dfrac{\pi}{2}\)，所以 \(\varphi=\dfrac{\pi}{3}\). 因此选 D.
\end{solution}
\end{problem}

\begin{problem}
直线 \(x - 2y + 1 = 0\) 关于直线 \(x = 1\) 对称的直线方程是（\quad）

\choices
    {\(x + 2y - 1 = 0\)}
    {\(2x + y - 1 = 0\)}
    {\(2x + y - 3 = 0\)}
    {\(x + 2y - 3 = 0\)}
\begin{answer}
D
\end{answer}

\begin{solution}
设原直线上的点为 \((u,v)\)，关于直线 \(x=1\) 对称后变为 \((x,y)=(2-u,v)\)，即 \(u=2-x\)，\(v=y\). 代入原直线方程得 \((2-x)-2y+1=0\)，即 \(x+2y-3=0\). 因此选 D.
\end{solution}
\end{problem}

\begin{problem}
要在边长为 \(16\text{ 米}\) 的正方形草坪上安装喷水龙头，使整个草坪都能喷洒到水. 假设每个喷水龙头的喷洒范围都是半径为 \(6\text{ 米}\) 的圆面，则需安装这种喷水龙头的个数最少是（\quad）

\bitmapfigure[max width=0.44\linewidth]{img/2007/zhejiang_science/q4_fig1.png}

\choices
    {\(3\)}
    {\(4\)}
    {\(5\)}
    {\(6\)}
\begin{answer}
B
\end{answer}

\begin{solution}
正方形的相邻两个顶点距离为 \(16\text{ 米}\)，而一个喷水龙头的喷洒直径为 \(12\text{ 米}\)，所以一个喷水龙头不可能同时覆盖两个顶点，至少需要 \(4\) 个喷水龙头.

把正方形分成 \(4\) 个边长为 \(8\text{ 米}\) 的小正方形，在每个小正方形的中心安装一个喷水龙头. 小正方形中任一点到其中心的距离不超过 \(4\sqrt{2}<6\text{ 米}\)，因此这 \(4\) 个喷水龙头可以覆盖整个草坪. 故最少安装 \(4\) 个，选 B.
\end{solution}
\end{problem}

\begin{problem}
已知随机变量 \(\xi\) 服从正态分布 \(N(2, \sigma^2)\)，\(P(\xi \leqslant 4) = 0.84\)，则 \(P(\xi \leqslant 0) =\)（\quad）

\choices
    {\(0.16\)}
    {\(0.32\)}
    {\(0.68\)}
    {\(0.84\)}
\begin{answer}
A
\end{answer}

\begin{solution}
正态分布的图象关于均值 \(2\) 对称，数值 \(0\) 与 \(4\) 到均值的距离相等且在两侧，故
\[
P(\xi\leqslant0)=P(\xi\geqslant4)=1-P(\xi\leqslant4)=1-0.84=0.16.
\]
因此选 A.
\end{solution}
\end{problem}

\begin{problem}
若 \(P\) 是两条异面直线 \(l\)，\(m\) 外的任意一点，则（\quad）

\choices
    {过点 \(P\) 有且仅有一条直线与 \(l\)，\(m\) 都平行}
    {过点 \(P\) 有且仅有一条直线与 \(l\)，\(m\) 都垂直}
    {过点 \(P\) 有且仅有一条直线与 \(l\)，\(m\) 都相交}
    {过点 \(P\) 有且仅有一条直线与 \(l\)，\(m\) 都异面}
\begin{answer}
B
\end{answer}

\begin{solution}
设两条异面直线 \(l\)，\(m\) 的方向向量分别为 \(\bs{u}\)，\(\bs{v}\). 由于 \(l\)，\(m\) 不平行，\(\bs{u}\times\bs{v}\ne\bs{0}\). 过点 \(P\) 且方向为 \(\bs{u}\times\bs{v}\) 的直线，方向同时垂直于 \(\bs{u}\) 和 \(\bs{v}\)，因而与 \(l\)，\(m\) 都垂直；其方向被唯一确定，所以这样的直线有且仅有一条. 因此选 B.
\end{solution}
\end{problem}

\begin{problem}
若非零向量 \(\bs{a}\)，\(\bs{b}\) 满足 \(|\bs{a} + \bs{b}| = |\bs{b}|\)，则（\quad）

\choices
    {\(|2\bs{a}| > |2\bs{a} + \bs{b}|\)}
    {\(|2\bs{a}| < |2\bs{a} + \bs{b}|\)}
    {\(|2\bs{b}| > |\bs{a} + 2\bs{b}|\)}
    {\(|2\bs{b}| < |\bs{a} + 2\bs{b}|\)}
\begin{answer}
C
\end{answer}

\begin{solution}
由 \(\lvert\bs{a}+\bs{b}\rvert=\lvert\bs{b}\rvert\)，平方得
\[
\lvert\bs{a}\rvert^2+2\bs{a}\cdot\bs{b}+\lvert\bs{b}\rvert^2=\lvert\bs{b}\rvert^2,
\]
所以 \(\bs{a}\cdot\bs{b}=-\dfrac12\lvert\bs{a}\rvert^2\). 于是
\[
\lvert\bs{a}+2\bs{b}\rvert^2
=\lvert\bs{a}\rvert^2+4\bs{a}\cdot\bs{b}+4\lvert\bs{b}\rvert^2
=4\lvert\bs{b}\rvert^2-\lvert\bs{a}\rvert^2
<4\lvert\bs{b}\rvert^2=\lvert2\bs{b}\rvert^2.
\]
故 \(\lvert2\bs{b}\rvert>\lvert\bs{a}+2\bs{b}\rvert\)，选 C.
\end{solution}
\end{problem}

\begin{problem}
设 \(f'(x)\) 是函数 \(f(x)\) 的导函数，将 \(y = f(x)\) 和 \(y = f'(x)\) 的图象画在同一个直角坐标系中，不可能正确的是（\quad）

\choices
    {\choicebitmap[width=0.15\paperwidth]{img/2007/zhejiang_science/q8_optA.png}}
    {\choicebitmap[width=0.15\paperwidth]{img/2007/zhejiang_science/q8_optB.png}}
    {\choicebitmap[width=0.15\paperwidth]{img/2007/zhejiang_science/q8_optC.png}}
    {\choicebitmap[width=0.15\paperwidth]{img/2007/zhejiang_science/q8_optD.png}}
\begin{answer}
D
\end{answer}

\begin{solution}
若上方曲线为 \(y=f(x)\)，则它在横坐标为正的位置有一个极小值点. 由导数与单调性的关系，在该点应有 \(f'(x)=0\)，所以下方曲线 \(y=f'(x)\) 必须经过横轴上的同一个点.

选项 D 中，上方曲线的极小值点横坐标为正，而下方曲线在该横坐标处位于横轴下方，不经过横轴，故不可能是同一函数 \(f\) 与其导函数 \(f'\) 的图象. 选 D.
\end{solution}
\end{problem}

\begin{problem}
已知双曲线 \(\frac{x^2}{a^2} - \frac{y^2}{b^2} = 1\)（\(a > 0\)，\(b > 0\)）的左，右焦点分别为 \(F_1\)，\(F_2\)，\(P\) 是准线上一点，且 \(PF_1 \perp PF_2\)，\(|PF_1| \cdot |PF_2| = 4ab\)，则双曲线的离心率是（\quad）

\choices
    {\(\sqrt{2}\)}
    {\(\sqrt{3}\)}
    {\(2\)}
    {\(3\)}
\begin{answer}
B
\end{answer}

\begin{solution}
记双曲线的半焦距为 \(c\)，则 \(c^2=a^2+b^2\). 不妨取右准线上的点 \(P\)，设其坐标为 \(\left(\dfrac{a^2}{c},y\right)\)，焦点为 \(F_1=(-c,0)\)，\(F_2=(c,0)\). 由 \(PF_1\perp PF_2\)，得
\[
\left(\dfrac{a^2}{c}+c,y\right)\cdot\left(\dfrac{a^2}{c}-c,y\right)=0,
\]
即 \(\left(\dfrac{a^2}{c}\right)^2+y^2=c^2\). 因此 \(P\) 在以原点为圆心、\(c\) 为半径的圆上.

由 \(PF_1\perp PF_2\)，还可得
\[
\lvert PF_1\rvert^2\lvert PF_2\rvert^2
=\bigl(2c^2+2c\cdot\tfrac{a^2}{c}\bigr)\bigl(2c^2-2c\cdot\tfrac{a^2}{c}\bigr)
=4c^2\left(c^2-\tfrac{a^4}{c^2}\right)=4c^2y^2.
\]
所以 \(\lvert PF_1\rvert\lvert PF_2\rvert=2c\lvert y\rvert\). 又
\[
y^2=c^2-\frac{a^4}{c^2}=\frac{(c^2-a^2)(c^2+a^2)}{c^2}=\frac{b^2(c^2+a^2)}{c^2},
\]
从而 \(\lvert PF_1\rvert\lvert PF_2\rvert=2b\sqrt{c^2+a^2}\). 根据题设 \(\lvert PF_1\rvert\lvert PF_2\rvert=4ab\)，得到 \(\sqrt{c^2+a^2}=2a\)，即 \(c^2=3a^2\). 故离心率 \(e=\dfrac ca=\sqrt3\)，选 B.
\end{solution}
\end{problem}

\begin{problem}
设 \( f(x) = \begin{cases}
x^2, & |x| \geqslant 1,\\
x, & |x| < 1,
\end{cases} \) 若 \( f(g(x)) \) 的值域是 \([0, +\infty)\)，则函数 \( g(x) \) 的值域是（\quad）

\choices
    {\((-\infty, -1] \cup [1, +\infty)\)}
    {\((-\infty, -1] \cup [0, +\infty)\)}
    {\([0, +\infty)\)}
    {\([1, +\infty)\)}
\begin{answer}
条件不足，无法唯一确定
\end{answer}

\begin{solution}
令 \(G\) 为函数 \(g(x)\) 的值域. 由
\[
f(u)=\begin{cases}
u^2,&\lvert u\rvert\geqslant1,\\
u,&\lvert u\rvert<1,
\end{cases}
\]
可知当 \(-1<u<0\) 时 \(f(u)<0\)，而 \(f(u)\geqslant0\) 的必要条件是 \(u\) 不属于 \((-1,0)\). 另一方面，为得到复合函数值域中的每个 \(y\in[0,1)\)，必须取到 \(u=y\in[0,1)\)；对于 \(y\geqslant1\)，既可以取 \(u=\sqrt y\)，也可以取 \(u=-\sqrt y\).

因此，仅由 \(f(g(x))\) 的值域为 \([0,+\infty)\)，不能唯一确定 \(G\). 例如函数 \(g_1(x)=\lvert x\rvert\) 的值域为 \(G_1=[0,+\infty)\)，且 \(f(G_1)=[0,+\infty)\). 再定义
\[
g_2(x)=\begin{cases}
x, & x\leqslant-1,\\
0, & -1<x<0,\\
x, & x\geqslant0,
\end{cases}
\]
则其值域为 \(G_2=(-\infty,-1]\cup[0,+\infty)\)，同样有 \(f(G_2)=[0,+\infty)\). 所以按原题面条件，函数 \(g(x)\) 的值域不唯一；若另加 \(G\) 为区间（例如 \(g\) 连续）的条件，才可得到 \(G=[0,+\infty)\)，对应选项 C.
\end{solution}
\end{problem}

\section{填空题}

\begin{problem}
已知复数 \(z_{1} = 1 - \i\)，\(z_{1} \cdot z_{2} = 1 + \i\)，则复数 \( z_{2} = \)\fillinblank{}.
\begin{answer}
\(\i\)
\end{answer}

\begin{solution}
由 \(z_1z_2=1+\i\)，得
\[
z_2=\frac{1+\i}{1-\i}=\frac{(1+\i)^2}{(1-\i)(1+\i)}=\frac{2\i}{2}=\i.
\]
\end{solution}
\end{problem}

\begin{problem}
已知 \(\sin \theta + \cos \theta = \frac{1}{5}\)，且 \(\frac{\pi}{2} \leqslant \theta \leqslant \frac{3\pi}{4}\)，则 \(\cos 2\theta\) 的值是\fillinblank{}.
\begin{answer}
\(-\frac{7}{25}\)
\end{answer}

\begin{solution}
由 \(\sin\theta+\cos\theta=\dfrac15\)，得
\[
\sin2\theta=\left(\sin\theta+\cos\theta\right)^2-1=-\frac{24}{25}.
\]
因为 \(\dfrac\pi2\leqslant\theta\leqslant\dfrac{3\pi}{4}\)，所以 \(\pi\leqslant2\theta\leqslant\dfrac{3\pi}{2}\). 又因 \(\sin2\theta=-\dfrac{24}{25}\ne0\)，故 \(\cos2\theta<0\). 再由
\[
\cos^22\theta=1-\sin^22\theta=1-\frac{576}{625}=\frac{49}{625},
\]
故 \(\cos2\theta=-\dfrac7{25}\).
\end{solution}
\end{problem}

\begin{problem}
不等式 \(\left|2x - 1\right| - x < 1\) 的解集是\fillinblank{}.
\begin{answer}
\(0 < x < 2\)
\end{answer}

\begin{solution}
当 \(x\geqslant\dfrac12\) 时，
\[
\lvert2x-1\rvert-x<1\Longleftrightarrow x-1<1\Longleftrightarrow x<2,
\]
故得到 \(\dfrac12\leqslant x<2\). 当 \(x<\dfrac12\) 时，
\[
\lvert2x-1\rvert-x<1\Longleftrightarrow1-3x<1\Longleftrightarrow x>0,
\]
故得到 \(0<x<\dfrac12\). 合并得解集为 \(0<x<2\).
\end{solution}
\end{problem}

\begin{problem}
某书店有 \(11\) 种杂志，\(2\text{ 元}\) \(1\text{ 本}\) 的 \(8\) 种，\(1\text{ 元}\) \(1\text{ 本}\) 的 \(3\) 种. 小张用 \(10\text{ 元}\) 钱买杂志（每种至多买一本，\(10\text{ 元}\) 钱刚好用完），则不同买法的种数是\fillinblank{}.（用数字作答）
\begin{answer}
\(266\)
\end{answer}

\begin{solution}
设购买 \(2\text{ 元}\) 杂志 \(i\) 种，购买 \(1\text{ 元}\) 杂志 \(j\) 种，则
\[
2i+j=10,\qquad0\leqslant i\leqslant8,\qquad0\leqslant j\leqslant3.
\]
所以只有 \((i,j)=(5,0)\) 或 \((4,2)\) 两种数量组合. 相应买法数为
\[
\mathrm C_8^5\mathrm C_3^0+\mathrm C_8^4\mathrm C_3^2=56+70\times3=266.
\]
故不同买法的种数为 \(266\).
\end{solution}
\end{problem}

\begin{problem}
随机变量 \(\xi\) 的分布列如下：

\begin{center}
\begin{tabular}{c|ccc}
\(\xi\) & \(-1\) & \(0\) & \(1\)\\ \hline
\(P\) & \(a\) & \(b\) & \(c\)
\end{tabular}
\end{center}

其中 \(a\)，\(b\)，\(c\) 成等差数列. 若 \(E(\xi) = \frac{1}{3}\)，则 \(D(\xi)\) 的值是\fillinblank{}.
\begin{answer}
\(\frac{5}{9}\)
\end{answer}

\begin{solution}
因为 \(a\)，\(b\)，\(c\) 成等差数列，所以 \(a+c=2b\). 又由分布列的概率和为 \(1\)，得 \(a+b+c=1\)，从而 \(b=\dfrac13\)，\(a+c=\dfrac23\). 由期望条件
\[
E(\xi)=-a+c=\frac13,
\]
解得 \(a=\dfrac16\)，\(c=\dfrac12\). 于是
\[
E(\xi^2)=a+c=\frac23,\qquad D(\xi)=E(\xi^2)-E(\xi)^2=\frac23-\frac19=\frac59.
\]
\end{solution}
\end{problem}

\begin{problem}
已知点 \(O\) 在二面角 \(\alpha-AB-\beta\) 的棱上，点 \(P\) 在 \(\alpha\) 内，且 \(\angle POB = 45^{\circ}\). 若对于 \(\beta\) 内异于 \(O\) 的任意一点 \(Q\)，都有 \(\angle POQ \geqslant 45^{\circ}\)，则二面角 \(\alpha-AB-\beta\) 的大小是\fillinblank{}.
\begin{answer}
\(90^{\circ}\)
\end{answer}

\begin{solution}
设 \(\theta\) 为二面角的大小，取沿 \(OB\) 的单位向量为 \(\bs{e}_1\)，取平面 \(\alpha\) 内垂直于 \(AB\) 且指向 \(P\) 的单位向量为 \(\bs{e}_2\)，再取平面 \(\beta\) 内垂直于 \(AB\) 的单位向量为 \(\bs{e}_3\). 则
\[
\bs{e}_1\perp\bs{e}_2,\qquad \bs{e}_1\perp\bs{e}_3,\qquad \bs{e}_2\cdot\bs{e}_3=\cos\theta.
\]
由 \(\angle POB=45^{\circ}\)，设 \(OP=r\)，则
\[
\overrightarrow{OP}=\frac r{\sqrt2}\bs{e}_1+\frac r{\sqrt2}\bs{e}_2.
\]
平面 \(\beta\) 内任意向量都可以写成 \(s\bs{e}_1+t\bs{e}_3\). 因此，\(\overrightarrow{OP}\) 在平面 \(\beta\) 上的投影长度为
\[
\left\lvert\operatorname{proj}_{\beta}\overrightarrow{OP}\right\rvert
=\sqrt{\left(\frac r{\sqrt2}\right)^2+\left(\frac r{\sqrt2}\cos\theta\right)^2}
=r\sqrt{\frac{1+\cos^2\theta}{2}}.
\]
所以 \(\overrightarrow{OP}\) 与平面 \(\beta\) 内非零向量的夹角的最小值 \(\varphi\) 满足
\[
\cos\varphi=\sqrt{\frac{1+\cos^2\theta}{2}}.
\]
题设对任意 \(Q\in\beta\) 且 \(Q\ne O\) 都有 \(\angle POQ\geqslant45^{\circ}\)，故 \(\varphi\geqslant45^{\circ}\)，从而
\[
\sqrt{\frac{1+\cos^2\theta}{2}}\leqslant\frac1{\sqrt2}.
\]
于是 \(\cos^2\theta=0\)，所以 \(\theta=90^{\circ}\).
\end{solution}
\end{problem}

\begin{problem}
设 \(m\) 为实数，若 \(\left\{(x,y)\mid x - 2y + 5 \geqslant 0,\ 3 - x \geqslant 0,\ mx + y \geqslant 0\right\} \subseteq
\left\{(x,y)\mid x^{2} + y^{2} \leqslant 25\right\}\)，则 \(m\) 的取值范围是\fillinblank{}.
\begin{answer}
\(0 \leqslant m \leqslant \frac{4}{3}\)
\end{answer}

\begin{solution}
三条边界直线分别为 \(y=\dfrac{x+5}{2}\)，\(x=3\) 和 \(y=-mx\). 若 \(2m+1\leqslant0\)，则前两条不等式与 \(y\geqslant-mx\) 的交集向左无界，不可能包含在圆 \(x^2+y^2\leqslant25\) 内，故 \(m>-\dfrac12\). 此时区域为三角形，其三个顶点为
\[
A=(3,4),\qquad B=(3,-3m),\qquad C=\left(-\frac5{2m+1},\frac{5m}{2m+1}\right).
\]
圆盘是凸集，故三角形包含于圆盘的充要条件是三个顶点都在圆盘内.

点 \(A\) 在圆周上. 由点 \(B\) 在圆盘内，得
\[
9+9m^2\leqslant25\Longleftrightarrow-\frac43\leqslant m\leqslant\frac43.
\]
由点 \(C\) 在圆盘内，得
\[
\frac{25(1+m^2)}{(2m+1)^2}\leqslant25
\Longleftrightarrow m(3m+4)\geqslant0.
\]
结合 \(m>-\dfrac12\)，可知 \(3m+4>0\)，所以 \(m\geqslant0\). 综合得 \(0\leqslant m\leqslant\dfrac43\).
\end{solution}
\end{problem}

\section{解答题}

\begin{problem}
已知 \(\triangle ABC\) 的周长为 \(\sqrt{2} + 1\)，且 \(\sin A + \sin B = \sqrt{2} \sin C\).

\begin{enumerate}
\item 求边 \(AB\) 的长；
\item 若 \(\triangle ABC\) 的面积为 \(\frac{1}{6}\sin C\)，求角 \(C\) 的度数.
\end{enumerate}

\begin{solution}
\begin{enumerate}
\item 设 \(BC=a\)，\(CA=b\)，\(AB=c\). 由正弦定理，
\[
\sin A+\sin B=\sqrt{2}\sin C
\Longleftrightarrow a+b=\sqrt{2}c.
\]
又 \(a+b+c=\sqrt{2}+1\)，所以
\[
(\sqrt{2}+1)c=\sqrt{2}+1,
\]
从而 \(c=AB=1\).
\item 三角形面积为
\[
\frac12ab\sin C=\frac16\sin C.
\]
由于 \(C\) 是三角形内角，\(\sin C>0\)，故 \(ab=\dfrac13\). 由 \(a+b=\sqrt{2}\)，得
\[
a^2+b^2=(a+b)^2-2ab=2-\frac23=\frac43.
\]
在三角形 \(ABC\) 中应用余弦定理，
\[
1=c^2=a^2+b^2-2ab\cos C
=\frac43-\frac23\cos C,
\]
所以 \(\cos C=\dfrac12\). 因为 \(0<C<\pi\)，故 \(C=60^{\circ}\).
\end{enumerate}
\end{solution}
\end{problem}

\begin{problem}
在如图所示的几何体中，\(EA \perp\) 平面 \(ABC\)，\(DB \perp\) 平面 \(ABC\)，\(AC \perp BC\)，且 \(AC = BC = BD = 2AE\)，\(M\) 是 \(AB\) 的中点.

\begin{enumerate}
\item 求证：\(CM \perp EM\)；
\item 求 \(CM\) 与平面 \(CDE\) 所成的角.
\end{enumerate}

\bitmapfigure[max width=0.42\linewidth]{img/2007/zhejiang_science/q19_fig1.png}

\begin{solution}
\begin{enumerate}
\item 建立空间直角坐标系，取 \(C\) 为原点，令 \(CA\) 为 \(x\) 轴正方向，\(CB\) 为 \(y\) 轴正方向，垂直于平面 \(ABC\) 的方向为 \(z\) 轴正方向. 由题设可取
\[
C=(0,0,0),\quad A=(2,0,0),\quad B=(0,2,0),\quad E=(2,0,1),\quad D=(0,2,2).
\]
因为 \(M\) 是 \(AB\) 的中点，所以 \(M=(1,1,0)\). 于是
\[
\overrightarrow{CM}=(1,1,0),\qquad \overrightarrow{EM}=(-1,1,-1),
\]
从而
\[
\overrightarrow{CM}\cdot\overrightarrow{EM}=-1+1+0=0.
\]
故 \(CM\perp EM\).
\item 由
\[
\overrightarrow{CD}=(0,2,2),\qquad \overrightarrow{CE}=(2,0,1),
\]
可取平面 \(CDE\) 的一个法向量为
\[
\bs{n}=\overrightarrow{CD}\times\overrightarrow{CE}=(2,4,-4).
\]
设 \(CM\) 与平面 \(CDE\) 所成的角为 \(\theta\)，则
\[
\sin\theta
=\frac{\lvert\overrightarrow{CM}\cdot\bs{n}\rvert}{\lvert\overrightarrow{CM}\rvert\lvert\bs{n}\rvert}
=\frac{\lvert2+4\rvert}{\sqrt2\sqrt{2^2+4^2+(-4)^2}}
=\frac6{\sqrt2\sqrt{36}}
=\frac1{\sqrt2}.
\]
因此 \(\theta=45^{\circ}\).
\end{enumerate}
\end{solution}
\end{problem}

\begin{problem}
如图，直线 \(y = kx + b\) 与椭圆 \(\frac{x^2}{4} + y^2 = 1\) 交于 \(A\)，\(B\) 两点，记 \(\triangle AOB\) 的面积为 \(S\).

\begin{enumerate}
\item 求在 \(k = 0\)，\(0 < b < 1\) 的条件下，\(S\) 的最大值；
\item 当 \(\left|AB\right| = 2\)，\(S = 1\) 时，求直线 \(AB\) 的方程.
\end{enumerate}

\bitmapfigure[max width=0.60\linewidth]{img/2007/zhejiang_science/q20_fig1.png}

\begin{solution}
\begin{enumerate}
\item 当 \(k=0\) 时，直线为 \(y=b\). 它与椭圆交于
\[
x=\pm2\sqrt{1-b^2},
\]
所以
\[
S=\frac12\cdot4\sqrt{1-b^2}\cdot b=2b\sqrt{1-b^2}.
\]
令 \(u=b^2\)，则 \(0<u<1\)，且
\[
S^2=4u(1-u)\leqslant1.
\]
因此 \(S\) 的最大值为 \(1\)，当 \(u=\dfrac12\)，即 \(b^2=\dfrac12\) 时取得.
\item 设直线为 \(y=kx+b\)，即 \(kx-y+b=0\). 因为 \(\lvert AB\rvert=2\)，\(S=1\)，所以原点到直线 \(AB\) 的距离为
\[
\frac{2S}{\lvert AB\rvert}=1.
\]
故
\[
\frac{\lvert b\rvert}{\sqrt{1+k^2}}=1,\qquad b^2=1+k^2.
\]
把 \(y=kx+b\) 代入椭圆方程，交点的横坐标 \(x_1,x_2\) 满足
\[
(1+4k^2)x^2+8kbx+4b^2-4=0.
\]
其判别式为
\[
\Delta=16(1+4k^2-b^2),
\]
因而
\[
\lvert x_1-x_2\rvert
=\frac{4\sqrt{1+4k^2-b^2}}{1+4k^2}.
\]
由于两交点在斜率为 \(k\) 的直线上，
\[
\lvert AB\rvert^2=(1+k^2)\lvert x_1-x_2\rvert^2=4.
\]
代入上式及前式，令 \(u=k^2\)，得
\[
16(1+u)\frac{1+4u-(1+u)}{(1+4u)^2}=4
\Longleftrightarrow12u(1+u)=(1+4u)^2
\Longleftrightarrow(2u-1)^2=0.
\]
所以 \(k=\pm\dfrac{\sqrt2}{2}\)，且由前式得 \(\lvert b\rvert=\dfrac{\sqrt6}{2}\). 因此直线 \(AB\) 的方程为 \(y=\dfrac{\sqrt2}{2}x+\dfrac{\sqrt6}{2}\)，\(y=\dfrac{\sqrt2}{2}x-\dfrac{\sqrt6}{2}\)，\(y=-\dfrac{\sqrt2}{2}x+\dfrac{\sqrt6}{2}\) 或 \(y=-\dfrac{\sqrt2}{2}x-\dfrac{\sqrt6}{2}\).
\end{enumerate}
\end{solution}
\end{problem}

\begin{problem}
已知数列 \(\{a_{n}\}\) 中的相邻两项 \(a_{2k-1}\)，\(a_{2k}\) 是关于 \(x\) 的方程 \(x^{2} - (3k + 2^{k})x + 3k \cdot 2^{k} = 0\) 的两个根，且 \(a_{2k-1} \leqslant a_{2k}\)（\(k = 1, 2, 3, \dots\)）.

\begin{enumerate}
\item 求 \(a_{1}\)，\(a_{3}\)，\(a_{5}\)，\(a_{7}\)；
\item 求数列 \(\{a_{n}\}\) 的前 \(2n\) 项和 \(S_{2n}\)；
\item 记 \(f(n) = \frac{1}{2}\left(\frac{|\sin n|}{\sin n} + 3\right)\)，\(T_n = \frac{(-1)^{f(2)}}{a_1a_2} + \frac{(-1)^{f(3)}}{a_3a_4} + \frac{(-1)^{f(4)}}{a_5a_6} + \cdots + \frac{(-1)^{f(n+1)}}{a_{2n-1}a_{2n}}\)，求证：\(\frac{1}{6} \leqslant T_n \leqslant \frac{5}{24}\)（\(n \in \N^*\)）.
\end{enumerate}

\begin{solution}
\begin{enumerate}
\item 方程
\[
x^2-(3k+2^k)x+3k\cdot2^k=0
\]
可因式分解为
\[
(x-3k)(x-2^k)=0.
\]
所以 \(a_{2k-1}\)，\(a_{2k}\) 是 \(3k\)，\(2^k\) 中较小者和较大者. 依次取 \(k=1,2,3,4\)，得
\[
(a_1,a_2)=(2,3),\quad(a_3,a_4)=(4,6),\quad
(a_5,a_6)=(8,9),\quad(a_7,a_8)=(12,16).
\]
故 \(a_1=2\)，\(a_3=4\)，\(a_5=8\)，\(a_7=12\).
\item 每一对相邻项的和为
\[
a_{2k-1}+a_{2k}=3k+2^k.
\]
因此
\[
S_{2n}=\sum_{k=1}^n(3k+2^k)
=3\frac{n(n+1)}2+\left(2^{n+1}-2\right)
=\frac{3n^2+3n}{2}+2^{n+1}-2.
\]
\item 令
\[
\varepsilon_j=(-1)^{f(j)}.
\]
由于 \(j\) 为正整数且 \(\pi\) 为无理数，\(\sin j\ne0\)，所以 \(\varepsilon_j\) 有定义且 \(\lvert\varepsilon_j\rvert=1\). 又因为 \(\sin2>0\)，\(\sin3>0\)，\(\sin4<0\)，所以 \(\varepsilon_2=\varepsilon_3=1\)，\(\varepsilon_4=-1\). 又
\[
a_{2k-1}a_{2k}=3k\cdot2^k,
\]
故
\[
T_n=\sum_{k=1}^n\frac{\varepsilon_{k+1}}{3k2^k}.
\]
当 \(n=1\) 时，\(T_1=\dfrac16\)；当 \(n=2\) 时，\(T_2=\dfrac16+\dfrac1{24}=\dfrac5{24}\).
当 \(n\geqslant3\) 时，由 \(\varepsilon_4=-1\)，以及
\[
\sum_{k=4}^{\infty}\frac1{3k2^k}
<\frac1{12}\sum_{k=4}^{\infty}\frac1{2^k}
=\frac1{96}
<\frac1{72}
=\frac1{3\cdot3\cdot2^3},
\]
可得
\[
T_n-T_2
=-\frac1{72}+\sum_{k=4}^n\frac{\varepsilon_{k+1}}{3k2^k}
\leqslant-\frac1{72}+\sum_{k=4}^{\infty}\frac1{3k2^k}<0.
\]
所以 \(T_n<\dfrac5{24}\). 另一方面，对 \(n\geqslant2\)，
\[
T_n\geqslant T_2-\sum_{k=3}^{\infty}\frac1{3k2^k}
>\frac5{24}-\frac19\sum_{k=3}^{\infty}\frac1{2^k}
=\frac5{24}-\frac1{36}
=\frac{13}{72}>\frac16.
\]
结合 \(T_1=\dfrac16\) 和 \(T_2=\dfrac5{24}\)，得到
\[
\frac16\leqslant T_n\leqslant\frac5{24}.
\]
\end{enumerate}
\end{solution}
\end{problem}

\begin{problem}
设 \( f(x) = \frac{x^3}{3} \)，对任意实数 \(t\)，记 \(g_t(x) = t^{\frac{2}{3}}x - \frac{2}{3}t\).

\begin{enumerate}
\item 求函数 \(y = f(x) - g_8(x)\) 的单调区间；
\item 求证：
\begin{circlelist}
\item 当 \(x > 0\) 时，\(f(x) \geqslant g_t(x)\) 对任意实数 \(t\) 成立；
\item 有且仅有一个正实数 \(x_0\)，使得 \(g_8(x_0) \geqslant g_t(x_0)\) 对任意实数 \(t\) 成立.
\end{circlelist}
\end{enumerate}

\begin{solution}
\begin{enumerate}
\item 令 \(h(x)=f(x)-g_8(x)\). 因为 \(8^{2/3}=4\)，所以
\[
h(x)=\frac{x^3}{3}-4x+\frac{16}{3},\qquad h'(x)=x^2-4=(x-2)(x+2).
\]
当 \(x<-2\) 或 \(x>2\) 时，\(h'(x)>0\)；当 \(-2<x<2\) 时，\(h'(x)<0\). 因此，\(y=f(x)-g_8(x)\) 在 \((-\infty,-2]\) 和 \([2,+\infty)\) 上单调递增，在 \([-2,2]\) 上单调递减.
\item 题设中参数 \(t\) 取任意实数. 在实数范围内，若 \(t=-s\) 且 \(s>0\)，则
\[
(-s)^{2/3}=\left(\sqrt[3]{-s}\right)^2=s^{2/3}.
\]
\begin{enumerate}
\item 取 \(x=1\)，\(t=-1\)，则
\[
f(1)=\frac13,\qquad g_{-1}(1)=1+\frac23=\frac53,
\]
所以 \(f(1)<g_{-1}(1)\)，题面所述的“不论 \(t\) 为何实数”命题不成立.
\item 对任意正实数 \(x_0\)，令 \(t=-s\)，则
\[
g_{-s}(x_0)=s^{2/3}x_0+\frac23s\longrightarrow+\infty
\qquad(s\longrightarrow+\infty),
\]
而
\[
g_8(x_0)=4x_0-\frac{16}{3}
\]
是定值. 因此不存在正实数 \(x_0\)，使 \(g_8(x_0)\geqslant g_t(x_0)\) 对任意实数 \(t\) 成立.
\end{enumerate}
因此，按题设“任意实数 \(t\)”理解，第二问的两个待证命题均不成立；若题面原意是“任意正实数 \(t\)”，则需另行讨论.
\end{enumerate}
\end{solution}
\end{problem}
